\subsection{Injective Functions}
\label{subsec:injective-functions}

\begin{tcolorbox}[title={Definition --- Injective}]
\begin{definition}[Injective Function]
\label{def:alg-injective}
Let $f:A\to B$.
The function $f$ is \textbf{injective} if
\[
\forall a_1,a_2 \in A,\;
f(a_1)=f(a_2) \Rightarrow a_1=a_2.
\]
\end{definition}
\end{tcolorbox}

\begin{remark*}[Standard quantified statement]
\[
\operatorname{Injective}(f,A,B)
\Longleftrightarrow
f:A\to B
\land
\forall a_1,a_2\in A,\; f(a_1)=f(a_2)\Rightarrow a_1=a_2.
\]
\end{remark*}

\begin{remark*}[Definition predicate reading]
$f$ is injective exactly when equal outputs force equal inputs.
\end{remark*}

\begin{remark*}[Negated quantified statement]
\[
\neg\operatorname{Injective}(f,A,B)
\Longleftrightarrow
\exists a_1,a_2\in A,\;
a_1\neq a_2 \land f(a_1)=f(a_2).
\]
\end{remark*}

\begin{remark*}[Interpretation]
Injectivity means distinct inputs do not collide under the function.
\end{remark*}

\begin{remark*}[Dependencies]
\begin{itemize}
  \item Function.
\end{itemize}
\end{remark*}
